%--- E11: Section 6.2, after the proof of Theorem 6.1 (before the concrete case)
\begin{explained}{the additive interaction remainder and why the terms commute}
\textbf{Step 1: the setup.} $\partial_tu=Au$, $u(0)=u_0$, with $A=\sum_{n=1}^NA_n$. The solution is formally $u(t)=e^{tA}u_0$. If term $n$ acted alone it would produce $S_n(t)=e^{tA_n}$, with change $\Delta_nu_0=(S_n(t)-\Id)u_0$ ($\Delta_n$ is the \emph{increment operator}). The additive interaction remainder is
\[
\Eint^{\mathrm{add}}(t)=u(t)-u_0-\sum_{n=1}^N\Delta_nu_0,
\]
that is, \emph{interaction $=$ full evolution $-$ initial state $-$ sum of the individual increments.}
Note that this is an additive decomposition, \emph{not} a decomposition into kernel spaces as in Section~\ref{sec:fact}.

\medskip\textbf{Step 2: why commutation gives an exact formula.} If $A_iA_j=A_jA_i$ then (for suitable operators and a well-defined evolution) $e^{t\sum_nA_n}=\prod_ne^{tA_n}=\prod_nS_n(t)$. Since $S_n=\Id+\Delta_n$,
\[
u(t)=\prod_{n=1}^N(\Id+\Delta_n)u_0,\qquad
\Eint^{\mathrm{add}}(t)=\Big[\prod_{n=1}^N(\Id+\Delta_n)-\Id-\sum_{n=1}^N\Delta_n\Big]u_0 .
\]
This is pure algebra: multiply the increments, then subtract the constant term and all terms involving a single increment.

\medskip\textbf{Step 3: $N=2$.} $u(t)=S_1S_2u_0=(\Id+\Delta_1)(\Id+\Delta_2)u_0=(\Id+\Delta_1+\Delta_2+\Delta_1\Delta_2)u_0$. Subtracting $u_0$ and the two single increments leaves
\[
\Eint^{\mathrm{add}}(t)=\Delta_1\Delta_2u_0=(S_1-\Id)(S_2-\Id)u_0 .
\]
It is an interaction term because it contains the effect of both increments acting together; if either increment is zero on the relevant state, the interaction vanishes.

\medskip\textbf{Step 4: $N=3$ shows the general structure.}
\[
(\Id+\Delta_1)(\Id+\Delta_2)(\Id+\Delta_3)=\Id+\Delta_1+\Delta_2+\Delta_3+\Delta_1\Delta_2+\Delta_1\Delta_3+\Delta_2\Delta_3+\Delta_1\Delta_2\Delta_3,
\]
so $\Eint^{\mathrm{add}}=(\Delta_1\Delta_2+\Delta_1\Delta_3+\Delta_2\Delta_3+\Delta_1\Delta_2\Delta_3)u_0$: three \emph{pairwise} interactions and one \emph{three-way} interaction. In general $\Eint^{\mathrm{add}}=\sum_{J\subseteq\{1,\dots,N\},|J|\ge2}\big(\prod_{j\in J}\Delta_j\big)u_0$: the empty subset gives the initial state, subsets of size one the individual increments, everything left involves at least two operators.

\medskip\textbf{Step 5: the Fourier formula.} If each $A_n$ is a constant-coefficient operator it becomes multiplication by its symbol $a_n(\kappa)$, so $\widehat{S_n(t)u_0}=e^{ta_n(\kappa)}\hat u_0$ and $\widehat{\Delta_nu_0}=(e^{ta_n(\kappa)}-1)\hat u_0$. Then
\[
\widehat{\Eint^{\mathrm{add}}}(t,\kappa)=\Big[\prod_ne^{ta_n(\kappa)}-1-\sum_n\big(e^{ta_n(\kappa)}-1\big)\Big]\hat u_0(\kappa).
\]
At each frequency this is ordinary scalar algebra; the operators commute because their Fourier symbols are multiplications. \emph{Notation note:} here $\kappa$ is the real frequency and $a_n(\kappa)$ the symbol of $A_n$ in the variable $i\kappa$ (e.g.\ $a(\kappa)=-ic\kappa-D\kappa^2$ for advection--diffusion).

\medskip\textbf{Step 6: small-time reading.} $\Delta_n=e^{tA_n}-\Id=tA_n+O(t^2)$. For two operators, $\Eint^{\mathrm{add}}=\Delta_1\Delta_2u_0=t^2A_1A_2u_0+O(t^3)$ (given enough regularity). Each single increment starts at order $t$, their interaction at order $t^2$: it needs both effects to contribute. For $N$ operators the interaction involving all $N$ increments starts at $t^N$, while the full remainder also contains pairwise interactions, generally starting at $t^2$.

\medskip\textbf{Step 7: why the terms commute.} Two operators that are polynomials in the \emph{same} $\partial_x$, e.g.\ $A_1=a_0+a_1\partial_x+a_2\partial_x^2$ and $A_2=b_0+b_1\partial_x+b_2\partial_x^2$, commute because powers of $\partial_x$ commute: $[A_1,A_2]=0$. More generally any two constant-coefficient operators, polynomials in $\partial_{x_1},\dots,\partial_{x_d}$, commute; for instance $A_1=\partial_x^2+\partial_y$ and $A_2=3\partial_x\partial_y-\partial_y^2$, because partial derivatives commute on sufficiently smooth functions (and on distributions). So for constant coefficients commutation is not a restrictive hypothesis: it comes almost automatically from the structure of the operators.

\medskip\textbf{The conceptual contrast.} Section~\ref{sec:fact} asked whether a solution can be stably separated into pieces belonging to different kernels. Here the question is different: if the generator is a sum of operators, how much of the resulting evolution cannot be explained by adding their individual changes? When the operators commute the answer is exact: expand the product of the evolution operators, remove the initial state and the single-operator contributions, and what remains is exactly the mixed interaction terms.
\end{explained}

